\documentclass[english,hebrew]{article}
\usepackage[HE8,T1]{fontenc}
\usepackage[utf8]{inputenc}
\setcounter{secnumdepth}{-2}
\setlength{\parindent}{0bp}
\usepackage{calc}
\usepackage{amsmath}
\usepackage{amsthm}
\usepackage{amssymb}

\makeatletter
\@ifundefined{date}{}{\date{}}
\usepackage{graphicx}
\usepackage{geometry}
\newcommand{\limi}{\underset{n\rightarrow\infty}{\lim}}
\newcommand{\limxz}{\underset{x\rightarrow x_{0}}{\lim}}
\newcommand{\limfxz}{\underset{x\rightarrow x_{0}}{\lim}f(x)}
\newcommand{\limfxm}{\underset{x\rightarrow x_{0}^{-}}{\lim}f(x)}
\newcommand{\limfxp}{\underset{x\rightarrow x_{0}^{+}}{\lim}f(x)}
\newcommand{\limxm}{\underset{x\rightarrow x_{0}^{-}}{\lim}}
\newcommand{\limxp}{\underset{x\rightarrow x_{0}^{+}}{\lim}}
\newcommand{\sqrm}{M_{m\times m}(\mathbb{F})}
\newcommand{\seqa}{(a_{n})_{n=1}^{\infty}}
\newcommand{\seqx}{(x_{n})_{n=1}^{\infty}}
\newcommand{\funcdr}{f:D\rightarrow\mathbb{R}}
\newcommand{\der}{\limxz\frac{f(x)-f(x_{0})}{x-x_{0}}}
\usepackage[all]{pdfprivacy}

\makeatother

\usepackage{babel}
\begin{document}
\newgeometry{top=1cm, bottom=2cm, left=1cm, right=1cm}
\title{אלגברה ליניארית {\beginL 2\endL} - הרצאה {\beginL 3\endL}}
\author{נעם לב}
\date{\L{$18.10.2026$}}

\maketitle
\noindent\fbox{\begin{minipage}[t]{1\columnwidth - 2\fboxsep - 2\fboxrule}%
\begin{center}
האוניברסיטה העברית \textbar{} אלגברה ליניארית ){\beginL 2\endL}( -
\L{$80135$} \textbar{} סמסטר{\small{} }{\beginL 2027\endL}{\small{} }א'
)תשפ\textquotedblleft ז א'( \textbar{} מרצה: פרופסור יבגני סטרחוב
\textbar{} מתרגל: מר לזר גוטרמן
\par\end{center}%
\end{minipage}}
\begin{center}
הסיכום הוא לא רשמי ולא עבר אישור של הסגל, על כן אין להשתמש בו )ובין
השאר במספור שלו ובשמות הטענות שבו( כמקור רשמי. אם ישנן טעויות, הן
שלי ולא של הסגל.
\par\end{center}

\medskip{}

\tableofcontents{}

\newpage{}

\section{משפט - תכונות אופרטור עם \L{$n$} ערכים עצמיים שונים}

יהי \L{$V$} מ\textquotedblright ו מעל שדה כלשהו \L{$\mathbb{F}$}
כך ש- \L{%
\[
\dim V=n\geq1
\]
}%
\L{%
\[
T\in\mathcal{L}(V)
\]
}%
ויש ל-\L{$T$} בדיוק \L{$n$} ערכים עצמיים \L{$\lambda_{1},\dots,\lambda_{n}$}
שונים זה מזה. אז:\\

)I( קיים בסיס \L{$\mathcal{B}_{V}$} של \L{$V$} כך שמתקיים \L{%
\[
M(T,\mathcal{B}_{V},\mathcal{B}_{V})=\begin{pmatrix}\lambda_{1} & 0_{\mathbb{F}} & \cdots & 0_{\mathbb{F}}\\
0_{\mathbb{F}} & \lambda_{2} & \cdots & 0_{\mathbb{F}}\\
\vdots &  & \ddots & \vdots\\
0_{\mathbb{F}} & 0_{\mathbb{F}} & \cdots & \lambda_{n}
\end{pmatrix}
\]
}%
כלומר, המטריצה המייצגת \L{$\left[T\right]_{\mathcal{B}}^{\mathcal{B}}$}
היא מטריצה אלכסונית.

הערה: מטריצה מייצגת כזאת זה המצב הטוב ביותר שיכול להיות. למשל, חזקות
של מטריצה או פונקציות של מטריצה כמו אקספוננט דרך טור, קל לחשב כשהמטריצה
אלכסונית. מצד אחד, יש הרבה דרישות, אך מצד שני זה נכון מעל כל שדה והחלק
החשוב ביותר במשפט הזה הוא )I(.\\

)II( לכל \L{$j\in\left\{ 1,\dots,n\right\} $}, \L{%
\[
\text{Null}(T-\lambda_{j}I)=\left\{ v\in V\mid T(v)=\lambda_{j}v\right\} 
\]
}%
הוא תמ\textquotedblright ו שמור תחת \L{$T$}.

זה למעשה אוסף כל הוקטורים העצמיים של הערך העצמי \L{$\lambda_{j}$}
יחד עם \L{$0_{V}$}.

אצל \L{$\text{Axler}$} הוא נקרא המרחב העצמי \L{$E(\lambda_{j},T)$}.

\L{$I$} כאן הוא \L{$\text{Id}_{V}$}.\\

)III(\L{%
\[
V=\text{Null}(T-\lambda_{1}I)\oplus\dots\oplus\text{Null}(T-\lambda_{n}I)
\]
}%
דהיינו, המרחב מתפרק לסכום ישר של התמ\textquotedblright ו-ים שמשרים
הערכים העצמיים. \\

)V\textbf{I}( נניח כי\L{%
\begin{align*}
T(v_{1}) & =\lambda_{1}v_{1}\\
 & \vdots\\
T(v_{n}) & =\lambda_{n}v_{n}
\end{align*}
}%
וכי \L{$\forall1\le i\le n:v_{i}\ne0_{V}$}. אזי\L{%
\begin{align*}
\text{Null}(T-\lambda_{1}I) & =\text{Span}(v_{1})\\
 & \vdots\\
\text{Null}(T-\lambda_{n}I) & =\text{Span}(v_{n})
\end{align*}
}%

\begin{proof}
\textbf{)I(} 1. נסמן \L{$\mathcal{B}_{V}=(v_{1},\dots,v_{n})$}, כאשר
\L{$v_{j}$} וקטור עצמי השייך ל-\L{$\lambda_{j}$}.

לפי המשפט {]}וקטורים עצמיים השייכים לערכים עצמיים שונים הם בת\textquotedblright ל{[}
מהרצאה 2, \L{$(v_{1},\dots,v_{n})$} בלתי תלויה לינארית, ויש בה \L{$n=\dim V$}
וקטורים, לכן {]}לפי משפט מלינארית {\beginL 1\endL}{[} היא בסיס של
\L{$V$}.

2. נחשב את המטריצה המייצגת בבסיס הזה: \L{%
\begin{align*}
T(v_{1}) & =\lambda_{1}v_{1}+0_{\mathbb{F}}v_{2}+\dots+0_{\mathbb{F}}v_{n}\\
T(v_{2}) & =0_{\mathbb{F}}v_{1}+\lambda_{2}v_{2}+\dots+0_{\mathbb{F}}v_{n}\\
 & \vdots\\
T(v_{n}) & =0_{\mathbb{F}}v_{1}+0_{\mathbb{F}}v_{2}+\dots+\lambda_{n}v_{n}
\end{align*}
}%
\L{%
\[
\Rightarrow M(T,\mathcal{B}_{V},\mathcal{B}_{V})=\begin{pmatrix}\lambda_{1} & 0_{\mathbb{F}} & \cdots & 0_{\mathbb{F}}\\
0_{\mathbb{F}} & \lambda_{2} & \cdots & 0_{\mathbb{F}}\\
\vdots &  & \ddots & \vdots\\
0_{\mathbb{F}} & 0_{\mathbb{F}} & \cdots & \lambda_{n}
\end{pmatrix}
\]
}%
\end{proof}
\begin{proof}
\textbf{)II(} 1. יהי \L{$j\in\left\{ 1,\dots,n\right\} $}, נסתכל
על \L{$\text{Null}(T-\lambda_{j}I)$}.

\L{$T-\lambda_{j}I$} זה אופרטור )כצ\textquotedblright ל של אופרטורים(.

בהרצאה {\beginL 15\endL} באלגברה לינארית {\beginL 1\endL} הוכחנו שגרעין
של כל ה\textquotedblright ל הוא תמ\textquotedblright ו, לכן \L{$\text{Null}(T-\lambda_{j}I)$}
זה תמ\textquotedblright ו ונותר להוכיח שזה תמ\textquotedblright ו
שמור תחת \L{$T$}.

{\beginL 2\endL}. יהי \L{$u\in\text{Null}(T-\lambda_{j}I)$} ונוכיח
כי \L{$T(u)\in\text{Null}(T-\lambda_{j}I)$}. מהגדרת הגרעין \L{%
\[
(T-\lambda_{j}I)(u)=0_{V}\implies T(u)=\lambda_{j}u
\]
}%

{\beginL 3\endL}. \L{%
\[
(T-\lambda_{j}I)\left(T(u)\right)\underset{(2)}{=}(T-\lambda_{j}I)(\lambda_{j}u)=T\left(\lambda_{j}u\right)-\lambda_{j}^{2}u=\lambda_{j}\cdot T\left(u\right)-\lambda_{j}^{2}u\underset{(2)}{=}\lambda_{j}^{2}u-\lambda_{j}^{2}u=0_{V}
\]
}%
\L{%
\[
\Rightarrow T(u)\in\text{Null}(T-\lambda_{j}I)
\]
}%
כנדרש.
\end{proof}
\begin{proof}
\textbf{)III(} 1. צריך להראות כי כל וקטור \L{$v\in V$} ניתן להציג
באופן יחיד כ- \L{%
\[
v=u_{1}+\dots+u_{n}
\]
}%
כאשר\L{%
\begin{align*}
u_{1}\in & \text{\ Null}(T-\lambda_{1}I)\\
 & \vdots\\
u_{n}\in & \text{\ Null}(T-\lambda_{n}I)
\end{align*}
}%

2. הוכחת קיום: נשתמש בבסיס \L{$\mathcal{B}_{V}$} מההוכחה של )I(.
אז \L{%
\[
\exists\alpha_{1},\dots,\alpha_{n}\in\mathbb{F}\text{ s.t. }v=\alpha_{1}v_{1}+\dots+\alpha_{n}v_{n}
\]
}%
ונסמן\L{%
\begin{align*}
u_{1}= & \ \alpha_{1}v_{1}\\
 & \vdots\\
u_{n}= & \ \alpha_{n}v_{n}
\end{align*}
}%
נשים לב כי לכל \L{$1\le i\le n$} מתקיים\L{%
\[
(T-\lambda_{i}I)\left(u_{i}\right)=T(u_{i})-\lambda_{i}u_{i}=T(\alpha_{i}v_{i})-\lambda_{i}\alpha_{i}v_{i}
\]
}%
\L{%
\[
=\alpha_{i}\left(T(v_{i})-\lambda_{i}v_{i}\right)=\alpha_{i}\left(\lambda_{i}v_{i}-\lambda_{i}v_{i}\right)=0_{V}
\]
}%
\L{%
\[
\Rightarrow u_{i}\in\text{Null}(T-\lambda_{i}I)
\]
}%
כלומר, ההצגה קיימת. 

{\beginL 3\endL}. הוכחת יחידות: אם\L{%
\[
u_{1}+\dots+u_{n}=v=u_{1}'+\dots+u_{n}'
\]
}%

כך שלכל \L{$1\le i\le n$} מתקיים \L{$u_{i},u_{i}'\in\text{Null}(T-\lambda_{i}I)$},
אז נגדיר\L{%
\[
\forall1\le i\le n,\quad w_{i}=u_{i}-u_{i}'
\]
}%
בנוסף, \L{$\text{Null}(T-\lambda_{i}I)$} הוא תמ\textquotedblright ו
לכן סגור לחיבור וכפל בסקלר מכאן \L{$w_{i}\in\text{Null}(T-\lambda_{i}I)$}

ונקבל\L{%
\[
w_{1}+\dots+w_{n}=v-v=0_{V}
\]
}%
נניח בשלילה שקיימים \L{$w_{i}\ne0_{V}$}.

אז הם לא שווים ל-\L{$0_{V}$} וכל אחד מהם מקיים \L{$w_{i}\in\text{Null}(T-\lambda_{i}I)$},
לכן הם וקטורים עצמיים.

אז גם הסכום שלהם שווה ל-\L{$0_{V}$} )כי אלה שכבר שווים ל-\L{$0_{V}$}
לא תורמים(, אבל הם וקטורים עצמיים של ערכים עצמיים שונים ולכן לפי המשפט
שהוכחנו בסוף הרצאה {\beginL 2\endL} הם בת\textquotedblright ל, וזאת
סתירה כי כאן הם שווים ל-\L{$0_{V}$} בתור צ\textquotedblright ל לא
טריוויאלי )עם מקדמים {\beginL 1\endL}(.

{\beginL 4\endL}. לכן לכל \L{$i$}, \L{$w_{i}=0_{V}$} מכאן \L{$u_{i}=u_{i}'$}
וההצגה היא יחידה כנדרש.
\end{proof}
הערה מיבגני: דרך קלה לזכור מה זה סכום ישר היא לחשוב על \L{$\mathbb{R}^{3}$}:
הוא סכום ישר של ציר ה-\L{$x$}, ציר ה-\L{$y$} וציר ה-\L{$z$}\L{$;$}
כל וקטור ניתן להציג באופן יחיד כסכום של וקטורים על כל אחד מהצירים
האלו.
\begin{proof}
\textbf{)VI(} 1. יהי \L{$1\le i\le n$}. 

אז מההנחה \L{$T(v_{i})-\lambda_{i}v_{i}=0_{V}$} לכן \L{$v_{i}\in\text{Null}(T-\lambda_{i}I)$}.

אבל \L{$\text{Null}(T-\lambda_{i}I)$} הוא תמ\textquotedblright ו,
לכן לפי המשפט שלמדנו בלינארית {\beginL 1\endL} לפיו פרוש של וקטור
הוא התמ\textquotedblright ו המינימלי המכיל אותו, \L{$\text{Span}(v_{i})\subseteq\text{Null}(T-\lambda_{i}I)$}.
בנוסף, מתקיים \L{$\dim\left(\text{Span}(v_{i})\right)=1$}.

2. לכן \L{$\dim\left(\text{Null}(T-\lambda_{i}I)\right)\ge1$}.

{\beginL 3\endL}. מצד שני, לפי )III(, \L{%
\[
V=\text{Null}(T-\lambda_{1}I)\oplus\dots\oplus\text{Null}(T-\lambda_{n}I)
\]
}%

לכן לפי המשפט בלינארית {\beginL 1\endL} המציין כי המימד של סכום ישר
הוא סכום המימדים\L{%
\[
n=\underset{\geq1}{\underbrace{\dim(\text{Null}(T-\lambda_{1}I))}}+\dots+\underset{\geq1}{\underbrace{\dim(\text{Null}(T-\lambda_{n}I))}}
\]
}%

מכאן, חייב להתקיים\L{%
\[
\forall1\le i\le n:\dim\left(\text{Null}(T-\lambda_{i}I)\right)=1
\]
}%
 {\beginL 4\endL}. קיבלנו ש-\L{%
\[
\forall1\le i\le n:\dim\left(\text{Null}(T-\lambda_{i}I)\right)=\dim\left(\text{Span}(v_{i})\right)
\]
}%

אבל תמ\textquotedblright ו המוכל בתמ\textquotedblright ו מאותו מימד
שווה לו, לכן\L{%
\[
\forall1\le i\le n:\text{Null}(T-\lambda_{i}I)=\text{Span}(v_{i})
\]
}%
\end{proof}

\section{טענה - אם ל־\L{$T_{A}$} יש \L{$n$} ערכים עצמיים שונים זה מזה, אז
\L{$A$} דומה למטריצה האלכסונית של הערכים העצמיים, ומטריצת המעבר היא
זו שעמודותיה הן הוקטורים העצמיים}

תהא \L{$A$} מטריצה מעל שדה כלשהו \L{$\mathbb{F}$}:

\L{%
\[
A=\begin{pmatrix}a_{1,1} & \cdots & a_{1,n}\\
\vdots &  & \vdots\\
a_{n,1} & \cdots & a_{n,n}
\end{pmatrix}\in\text{Mat}(n\times n,\mathbb{F})
\]
}%
ונסתכל על המרחב \L{%
\[
V=\mathbb{F}^{n}=\left\{ \begin{pmatrix}d_{1}\\
\vdots\\
d_{n}
\end{pmatrix}\,\middle|\,\forall1\le i\le n:d_{i}\in\mathbb{F}\right\} 
\]
}%

נסמן ב-\L{$\check{v}$} איבר ב-\L{$\mathbb{F}^{n}$}, כלומר \L{$\check{v}\in\text{Mat}(n\times1,\mathbb{F})$}. 

בנוסף נגדיר\L{%
\[
T_{A}\in\mathcal{L}(\mathbb{F}^{n},\mathbb{F}^{n}):T_{A}(\check{v})\overset{\text{def}}{=}A\check{v}
\]
}%
)זו העתקה לינארית מ-\L{$\mathbb{F}^{n}$} ל-\L{$\mathbb{F}^{n}$},
כפל מטריצה \L{$n\times n$} במטריצה \L{$n\times1$}.(

נניח כי \L{%
\begin{align*}
T_{A}(\check{v}_{1}) & =A\check{v}_{1}=\lambda_{1}\check{v}_{1},\quad\lambda_{1}\in\mathbb{F}\\
 & \vdots\\
T_{A}(\check{v}_{n}) & =A\check{v}_{n}=\lambda_{n}\check{v}_{n},\quad\lambda_{n}\in\mathbb{F}
\end{align*}
}%
כאשר \L{$\check{v}_{1},\dots,\check{v}_{n}\neq\check{0}$} וכאשר \L{$\lambda_{1},\dots,\lambda_{n}$}
שונים זה מזה.\\

\emph{הערה על הטרמינולוגיה: \L{$\lambda_{1},\dots,\lambda_{n}$} נקראים
ערכים עצמיים של המטריצה \L{$A$}, ו-\L{$\check{v}_{1},\dots,\check{v}_{n}$}
נקראים וקטורים עצמיים של המטריצה \L{$A$}. יבגני: זו השפה של מהנדסים
ושל חלק מהספרות. הכוונה המתמטית היא לערכים ולוקטורים העצמיים של ההעתקה
\L{$T_{A}$}.}\\

אז עבור המטריצה שעמודותיה הן הוקטורים העצמיים:\L{%
\[
Q=\begin{pmatrix}\check{v}_{1} & \check{v}_{2} & \cdots & \check{v}_{n}\end{pmatrix}\in\text{Mat}(n\times n,\mathbb{F})
\]
}%

מתקיים

\L{%
\[
(*)\quad A=Q\begin{pmatrix}\lambda_{1} & 0_{\mathbb{F}} & \cdots & 0_{\mathbb{F}}\\
0_{\mathbb{F}} & \lambda_{2} & \cdots & 0_{\mathbb{F}}\\
\vdots &  & \ddots & \vdots\\
0_{\mathbb{F}} & 0_{\mathbb{F}} & \cdots & \lambda_{n}
\end{pmatrix}Q^{-1}
\]
}%

\begin{proof}
1. צריך להראות כי \L{$(*)$} מוגדר היטב, כלומר כי \L{$Q$} הפיכה.

\L{$\check{v}_{1},\dots,\check{v}_{n}$} הם וקטורים עצמיים השייכים
לערכים עצמיים שונים זה מזה של ההעתקה הלינארית \L{$T_{A}$}.

מכאן \L{$\check{v}_{1},\dots,\check{v}_{n}$} בלתי תלויים לינארית
לכן העמודות של \L{$Q$} בת\textquotedblright ל, אז \L{$Q$} הפיכה
ו-\L{$Q^{-1}$} קיימת.

2. \L{$(*)$} שקולה למשוואה \L{%
\[
(**)\quad AQ=Q\begin{pmatrix}\lambda_{1} & 0_{\mathbb{F}} & \cdots & 0_{\mathbb{F}}\\
0_{\mathbb{F}} & \lambda_{2} & \cdots & 0_{\mathbb{F}}\\
\vdots &  & \ddots & \vdots\\
0_{\mathbb{F}} & 0_{\mathbb{F}} & \cdots & \lambda_{n}
\end{pmatrix}
\]
}%
)מכפילים את שני האגפים ב-\L{$Q$} מימין, ו-\L{$Q^{-1}Q$} מצטמצם(

3. נחשב: \L{%
\begin{align*}
AQ & =A\begin{pmatrix}\check{v}_{1} & \dots & \check{v}_{n}\end{pmatrix}\underset{(1)}{=}\begin{pmatrix}A\check{v}_{1} & \dots & A\check{v}_{n}\end{pmatrix}=\begin{pmatrix}\lambda_{1}\check{v}_{1} & \dots & \lambda_{n}\check{v}_{n}\end{pmatrix}\\
 & =\underset{Q}{\underbrace{\begin{pmatrix}\check{v}_{1} & \cdots & \check{v}_{n}\end{pmatrix}}}\begin{pmatrix}\lambda_{1} & 0_{\mathbb{F}} & \dots & 0_{\mathbb{F}}\\
0_{\mathbb{F}} & \lambda_{2} & \dots & 0_{\mathbb{F}}\\
\vdots &  & \ddots & \vdots\\
0_{\mathbb{F}} & 0_{\mathbb{F}} & \dots & \lambda_{n}
\end{pmatrix}
\end{align*}
}%
\L{$(1)$} כפל מטריצה במטריצה הוא כפל בכל עמודה ועמודה.

קיבלנו בדיוק את \L{$(**)$}, כנדרש. 
\end{proof}
\L{$Q$} חייבת להיות משמאל. למה?

כי כפל במטריצה אלכסונית מימין מכפיל את העמודה ה-\L{$j$} ב-\L{$\lambda_{j}$},
ואילו כפל משמאל מכפיל את השורות. 

יבגני: תשכנעו את עצמכם שהחישוב נכון. אם אין מספיק ניסיון עם כפל מטריצות,
לעשות דוגמאות.\\

האלגוריתם: מוצאים ערכים עצמיים )בכל דרך(, בודקים שהם שונים זה מזה
ושמספרם הוא \L{$n$}, מוצאים וקטורים עצמיים מהערכים העצמיים השונים
ולכן מקבלים וקטורים עצמיים בת\textquotedblright ל, בונים מהם את \L{$Q$}
)כעמודות(, ואז \L{$A=QDQ^{-1}$} עם \L{$D$} אלכסונית.

ל-\L{$Q$} יש תפקיד גיאומטרי של מעבר בסיס, אבל כאן זה פשוט שוויון
מטריציוני.\\

הקשר ללינארית 1: \L{$A=M(T_{A},\varepsilon,\varepsilon)$} בבסיס הסטנדרטי
\L{$\varepsilon$}, \L{$Q=[\text{Id}]_{\varepsilon}^{\mathcal{B}}$}
היא מטריצת המעבר מהבסיס \L{$\mathcal{B}=(\check{v}_{1},\dots,\check{v}_{n})$}
לסטנדרטי, ו-\L{$(*)$} היא בדיוק \L{$M(T_{A},\mathcal{B},\mathcal{B})=Q^{-1}AQ$},
כלומר \L{$A$} דומה למטריצה אלכסונית.

\textquotedblright לכסון מטריצה\textquotedblright{} פירושו למצוא \L{$Q$}
כזאת.\\

זו תוצאה בסיסית בתורה הספקטרלית של העתקות לינאריות: במימד סופי, אם
יודעים מראש שיש \L{$n$} ערכים עצמיים שונים, ניתן ללכסן. 

במימד אינסופי אין תוצאות שנוכל להציג בקורס ולכן נישאר במימד סופי )הנושא
שייך לאנליזה פונקציונלית ודורש כלים מאנליזה. לדוגמה, המשפט הספקטרלי
לאופרטורים קומפקטיים צמודים לעצמם(.

השדה עדיין שדה כלשהו, אבל מעל \L{$\mathbb{C}$} יש משהו לומר אפילו
במקרה הכללי, בלי הנחה על הערכים העצמיים.

\section{הגדרה - מטריצה משולשית עליונה \L{$(\text{Upper triangular matrix})$}}

\L{%
\[
M=\begin{pmatrix}a_{1,1} & a_{1,2} & \cdots & a_{1,n}\\
0_{\mathbb{F}} & a_{2,2} & \cdots & a_{2,n}\\
\vdots & \ddots & \ddots & \vdots\\
0_{\mathbb{F}} & \cdots & 0_{\mathbb{F}} & a_{n,n}
\end{pmatrix}
\]
}%
כלומר, \L{$a_{i,j}=0_{\mathbb{F}}$} כאשר \L{$i>j$} )מספר השורה גדול
ממספר העמודה(. 

מטריצה מסוג זה נקראת מטריצה משולשית עליונה.

אין שום דרישה על האיברים על האלכסון ומעליו, כל סקלר מתאים.

בפרט, כל מטריצה \L{$1\times1$} היא משולשית עליונה.

\section{למה - אפיון מטריצה משולשית עליונה}

יהי \L{$V$} מ\textquotedblright ו מעל שדה כלשהו \L{$\mathbb{F}$}
כך ש- \L{%
\[
\dim V=n\geq1
\]
}%
\L{%
\[
T\in\mathcal{L}(V)
\]
}%
\L{%
\[
\mathcal{B}_{V}=(v_{1},\dots,v_{n})
\]
}%
אז התנאים הבאים שקולים:\\

)I( \L{$M(T,\mathcal{B}_{V},\mathcal{B}_{V})$} משולשית עליונה.

)II( לכל \L{$1\leq j\leq n$}, \L{$T(v_{j})\in\text{Span}(v_{1},\dots,v_{j})$}.

)III( לכל \L{$1\leq j\leq n$}, \L{$U_{j}=\text{Span}(v_{1},\dots,v_{j})$}
הוא תת-מרחב שמור תחת \L{$T$}.
\begin{proof}
\textbf{)I( \L{$\iff$} )II(}

נסמן \L{%
\[
M(T,\mathcal{B}_{V},\mathcal{B}_{V})=\begin{pmatrix}a_{1,1} & \cdots & a_{1,n}\\
\vdots &  & \vdots\\
a_{n,1} & \cdots & a_{n,n}
\end{pmatrix}
\]
}%
לפי הגדרת המטריצה המייצגת, \L{%
\[
T(v_{j})=\sum_{i=1}^{n}\left(a_{i,j}\cdot v_{i}\right)
\]
}%
לכן, לפי משפט יחידות ההצגה לפי בסיס,\L{%
\begin{align*}
T(v_{1})=\sum_{i=1}^{n}a_{i,1}v_{i}=a_{1,1}v_{1} & \iff T(v_{1})\in\text{Span}(v_{1})\\
T(v_{2})=\sum_{i=1}^{n}a_{i,2}v_{i}=a_{1,2}v_{1}+a_{2,2}v_{2} & \iff T(v_{2})\in\text{Span}(v_{1},v_{2})\\
 & \vdots\\
T(v_{j})=a_{1,j}v_{1}+\dots+a_{j,j}v_{j} & \iff T(v_{j})\in\text{Span}(v_{1},\dots,v_{j})
\end{align*}
}%
באגף שמאל: \L{$a_{i,j}=0_{\mathbb{F}}$} לכל \L{$i>j$}, כלומר המטריצה
משולשית עליונה: תנאי )I(.

באגף ימין: תנאי )II(.

\textbf{)II( \L{$\Leftarrow$} )III(}

1. יהי \L{$u\in\text{Span}(v_{1},\dots,v_{j})$}. אז קיימים \L{$d_{1},\dots,d_{j}$}
כך ש- \L{%
\[
u=d_{1}v_{1}+\dots+d_{j}v_{j}\implies T(u)=d_{1}T(v_{1})+\dots+d_{j}T(v_{j})
\]
}%

2. אבל לפי )II(, \L{%
\begin{align*}
T(v_{1})\in\text{Span}(v_{1}) & \subseteq\text{Span}(v_{1},\dots,v_{j})\\
 & \vdots\\
T(v_{j})\in\  & \text{Span}(v_{1},\dots,v_{j})
\end{align*}
}%
\L{$\Leftarrow$} \L{$T(u)\in\text{Span}(v_{1},\dots,v_{j})$}, כלומר
\L{$U_{j}$} שמור תחת \L{$T$}.

\textbf{)III( \L{$\Leftarrow$} )II(}

1. נניח כי \L{%
\[
\forall1\le j\le n:U_{j}=\text{Span}(v_{1},\dots,v_{j})
\]
}%
 הוא תמ\textquotedblright ו שמור תחת \L{$T$}.

יהי \L{$1\le j\le n$}. באופן טריוויאלי,\L{%
\[
v_{j}\in\text{Span}\left(v_{1},\dots,v_{j}\right)=U_{j}
\]
}%
 אבל \L{$U_{j}$} שמור תחת \L{$T$} לכן \L{$T(v_{j})\in U_{j}$}.

כלומר \L{$T(v_{j})\in\text{Span}\left(v_{1},\dots,v_{j}\right)$}
כנדרש.
\end{proof}

\section{משפט - מעל \L{$\mathbb{C}$}, כל אופרטור ניתן לייצוג על ידי מטריצה
משולשית עליונה}

יהי \L{$V$} מ\textquotedblright ו \textbf{מעל \L{$\mathbb{C}$}}
כך ש- \L{%
\[
\dim V=n\geq1
\]
}%
\L{%
\[
T\in\mathcal{L}(V)
\]
}%
אז קיים בסיס \L{$\mathcal{B}_{V}$} של \L{$V$} כך ש-\L{$M(T,\mathcal{B}_{V},\mathcal{B}_{V})$}
משולשית עליונה.\\

יבגני: זה אחד המשפטים הבסיסיים בתורה הספקטרלית, וכנראה ברשימה למבחן.\\

זה לא נכון מעל כל שדה, אפילו לא מעל \L{$\mathbb{R}$}. 

דוגמה נגדית אפשרית היא הסיבוב ב-\L{$90^{\circ}$} מהרצאה 2. 

למטריצה משולשית עליונה יש ערך עצמי )\L{$a_{1,1}$}, כי \L{$T(v_{1})=a_{1,1}v_{1}$}(,
ולסיבוב אין ערכים עצמיים ממשיים.
\begin{proof}
נוכיח באינדוקציה )חזקה( על \L{$n=\dim V$}.

1( עבור \L{$n=1$}:

\L{%
\[
\mathcal{B}_{V}=\left(v_{1}\right)
\]
}%
\L{%
\[
\therefore\exists a_{1,1}\in\mathbb{F}=\mathbb{C}\text{ s.t. }T(v_{1})=a_{1,1}v_{1}
\]
}%
\L{%
\[
M(T,\mathcal{B}_{V},\mathcal{B}_{V})=\left(a_{1,1}\right)
\]
}%
וכל מטריצה \L{$1\times1$} היא משולשית עליונה, לכן המשפט מתקיים עבור
\L{$n=1$}.

2( יהי \L{$n\geq2$}.

נניח כי המשפט מתקיים עבור כל מרחב וקטורי \L{$\widetilde{V}$} מעל
\L{$\mathbb{C}$} עם \L{$1\leq\dim\widetilde{V}\leq n-1$}, ועבור
כל אופרטור לינארי \L{$\widetilde{T}\in\mathcal{L}(\widetilde{V})$}.

נוכיח עבור \L{$V$} עם \L{$\dim V=n$}.

3. קיים ערך עצמי \L{$\lambda$} של \L{$T$} )כי \L{$V$} הוא מ\textquotedblright ו
נ\textquotedblright ס \textbf{מעל \L{$\mathbb{C}$}}, אז לפי המשפט
מהרצאה 2(.

4. נגדיר \L{%
\[
U\overset{\text{def}}{=}\text{Range}(T-\lambda I)
\]
}%
\L{$\lambda$} ערך עצמי, ולכן קיים \L{$v\neq0_{V}$} עם \L{$T(v)=\lambda v$},
כלומר \L{$(T-\lambda I)(v)=0_{V}$}.

\L{$0_{V}$} גם בגרעין, ולכן \L{$T-\lambda I$} אינה חח\textquotedblright ע.

מכאן, היא גם לא על, כי אופרטור על מרחב נ\textquotedblright ס הוא חח\textquotedblright ע
אם\textquotedblright ם הוא על.

אז\L{%
\[
\dim\text{Range}(T-\lambda I)\leq n-1
\]
}%
)אילו היה מתקיים \L{$\dim U=n$}, אז \L{$U=V$} וההעתקה הייתה על,
בסתירה(.

5. נסתכל על \L{$T|_{U}$} )הצמצום של \L{$T$} ל-\L{$U$}(.

א. כדי ש-\L{$T|_{U}$} יהיה אופרטור על \L{$U$} )כלומר \L{$T|_{U}\in\mathcal{L}(U)$}(
צריך ש-\L{$U$} יהיה שמור תחת \L{$T$}.

לכל \L{$u\in U$} מתקיים: \L{%
\[
T(u)=T\left(u\right)-\lambda u+\lambda u=(T-\lambda I)(u)+\lambda u
\]
}%
כאשר \L{%
\[
(T-\lambda I)(u)\in\text{Range}(T-\lambda I)=U
\]
}%
וגם \L{$\lambda u\in U$} כי \L{$U$} הוא תמ\textquotedblright ו.מכאן,
\L{$T(u)\in U$}.\\

ב. הנחת האינדוקציה דורשת גם \L{$\dim U\geq1$}.

אם זה לא המקרה, לא נשתמש בה, אלא בנימוק הבא:

אם \L{$\dim U=0$}, אז \L{$T-\lambda I=0$}, כלומר \L{$T=\lambda I$},
ואז בכל בסיס \L{$M(T,\mathcal{B}_{V},\mathcal{B}_{V})=\lambda I_{n}$},
שהיא אלכסונית ובפרט משולשית עליונה. במקרה הזה \L{$m=0$}.\\

ג. לפי הנחת האינדוקציה )2(, קיים בסיס \L{$\mathcal{B}_{U}=(u_{1},\dots,u_{m})$}
של \L{$U$} כך ש-\L{$M(T|_{U},\mathcal{B}_{U},\mathcal{B}_{U})$}
משולשית עליונה.

לפי למת האפיון, זה אומר כי \L{%
\[
\forall1\le j\le m:T(u_{j})\in\text{Span}(u_{1},\dots,u_{j})
\]
}%

6. נשלים את \L{$(u_{1},\dots,u_{m})$} )שהיא בת\textquotedblright ל
כבסיס של \L{$U$} ולכן אפשר להשלים אותה לבסיס של \L{$V$} לפי משפט
ההשלמה מהקורס לינארית {\beginL 1\endL}( לבסיס של \L{$V$}: \L{%
\[
\mathcal{B}_{V}=(\underset{n}{\underbrace{u_{1},\dots,u_{m},v_{1},\dots,v_{k}}})
\]
}%

כאשר \L{$k=n-m$}. לכל \L{$1\le j\le k$} מתקיים: \L{%
\[
T\left(v_{j}\right)=T\left(v_{j}\right)-\lambda v_{j}+\lambda v_{j}=\underset{\in\text{Range}(T-\lambda I)=U}{\underbrace{(T-\lambda I)(v_{j})}}+\lambda v_{j}
\]
}%
\L{%
\[
\implies T(v_{j})\in\text{Span}(u_{1},\dots,u_{m},v_{j})\subseteq\text{Span}(u_{1},\dots,u_{m},v_{1},\dots,v_{j})
\]
}%
יחד עם 5, כל וקטור בבסיס \L{$\mathcal{B}_{V}$} נשלח לספאן של הוקטורים
שלפניו )כולל עצמו(.

אז שוב לפי למת האפיון, \L{$M(T,\mathcal{B}_{V},\mathcal{B}_{V})$}
משולשית עליונה.
\end{proof}

\subsection{צמצום של העתקה לתת-מרחב}

אם \L{$T:V\rightarrow V$} העתקה לינארית ו-\L{$U\subseteq V$} תמ\textquotedblright ו
שמור, אפשר להתרכז רק בוקטורים שבתוך \L{$U$}.

במקרה זה, מתעסקים עם הצמצום הזה במקום עם ההעתקה המקורית, שמסומן \L{$T|_{U}$}
)בפועל פשוט שוכחים מכל הוקטורים שמחוץ ל-\L{$U$}(.\L{%
\[
T|_{U}:U\rightarrow U,\ T|_{U}(u)=T(u)
\]
}%

לעיתים לא מציינים זאת במפורש אלא אומרים \textquotedblright נתרכז בתמ\textquotedblright ו
\L{$U$} עם אותה \L{$T$}\textquotedblright .

זה מוגדר היטב )עם טווח \L{$U$}( בדיוק כאשר \L{$U$} שמור תחת \L{$T$},
וזו העתקה לינארית כי \L{$T$} לינארית.
\end{document}
